El recorrido realizado a lo largo de este capítulo pone de manifiesto que los problemas que motivaron el presente trabajo tienen raíces tanto pedagógicas como técnicas. Desde el ámbito de la enseñanza, la investigación evidencia que las dificultades en los cursos iniciales de programación (CS1) no responden únicamente a la complejidad intrínseca de los contenidos, sino también a la falta de retroalimentación personalizada en tiempo real y al riesgo del facilismo algorítmico que introduce el uso indiscriminado de herramientas de IA generativa. El método socrático emerge como un marco pedagógico sólido para paliar ambos problemas, puesto que, al sustituir la entrega directa de respuestas por un diálogo guiado de preguntas estratégicas, fomenta el andamiaje cognitivo y preserva el esfuerzo de razonamiento autónomo del estudiante. Desde la vertiente técnica, el análisis de los LLMs de código abierto demuestra que actualmente es factible construir un sistema de tutoría local, sin dependencia de APIs comerciales, capaz de soportar interacciones fluidas y continuas. Este panorama justifica el desarrollo de SocratiCode, una solución integral que combina ambos mundos para priorizar el descubrimiento guiado y la autonomía del alumno en la construcción de su propio pensamiento computacional.